% Generated draft: requires mathematical and source review.
\documentclass[11pt]{article}
\usepackage[margin=25mm]{geometry}
\usepackage{fontspec}
\setmainfont{Latin Modern Roman}
\usepackage{amsmath,amssymb,mathtools,unicode-math}
\setmathfont{Latin Modern Math}
\usepackage{enumitem,needspace,xurl,hyperref}
\setlength{\parindent}{0pt}
\setlength{\parskip}{5pt}
\setlength{\emergencystretch}{3em}
\newcommand{\N}{\mathbb{N}}
\newcommand{\Z}{\mathbb{Z}}
\newcommand{\Q}{\mathbb{Q}}
\newcommand{\R}{\mathbb{R}}
\newcommand{\C}{\mathbb{C}}
\newcommand{\F}{\mathbb{F}}
\newcommand{\A}{\mathbb{A}}
\newcommand{\PP}{\mathbb{P}}
\newcommand{\E}{\mathbb{E}}
\newcommand{\eps}{\varepsilon}
\newcommand{\dd}{\,\mathrm{d}}
\begin{document}
\section*{UnsolvedMath ID: 701426}
\textbf{Source problem code:} catalogue entry
\subsection{Definitions and mathematical statement}
For a finite group \(K\), define
\[
\pi(K)=\{p\in\N:p\text{ is prime and }p\mid |K|\}.
\]
In particular, \(\pi(\{1\})=\varnothing\). This is the standard interpretation of the notation \(\pi\), which the supplied record does not explicitly define. For \(a,b\in G\), let \(\langle a,b\rangle\) denote the subgroup generated by \(a\) and \(b\). The elements \(a\) and \(b\) are conjugate in \(G\) if \(b=x^{-1}ax\) for some \(x\in G\).

The question is whether every finite group \(G\) admits elements \(a,b,x\in G\) such that
\[
b=x^{-1}ax
\qquad\text{and}\qquad
\pi(G)=\pi(\langle a,b\rangle).
\]
No nonidentity or distinctness requirement is imposed on these elements.

The supplied source states that the assertion holds for soluble groups. Here soluble means that the derived series \(G^{(0)}=G\), \(G^{(i+1)}=[G^{(i)},G^{(i)}]\) reaches the trivial subgroup, where \([L,L]\) is the subgroup generated by the commutators \(u^{-1}v^{-1}uv\) with \(u,v\in L\).

The source's comment of 2013 reports that A. Lucchini, M. Morigi, and P. Shumyatsky proved that every finite group \(G\) has a subgroup \(H=\langle u,v\rangle\) satisfying \(\pi(G)=\pi(H)\). This reported result does not impose conjugacy on \(u\) and \(v\); the conjugacy requirement is part of the question above. These background claims are preserved as supplied and have not been independently verified.
\subsection{Short English statement}
Does every finite group contain two elements conjugate within that group whose generated subgroup has exactly the same prime divisors of its order as the whole group? The supplied record labels the problem open, states that it holds for soluble groups, and reports a result allowing two generators without a conjugacy requirement.
\subsection{Sources}
\begin{itemize}
\item Local UnsolvedMath registry: permanent ID \(701426\), exact source problem code \texttt{catalogue entry}, title \emph{Kourovka Notebook Problem 17.125}. Supplied upstream revision: \texttt{d371ee917ac9647fbb06cd6fc529deec6f3dda82}. The registry labels the problem open; this is supplied metadata, not an independently verified literature assessment.
\item \emph{The Kourovka Notebook (Unsolved Problems in Group Theory)}, Problem 17.125, proposed by P. Shumyatsky, 17th issue (2010). The supplied background also identifies the 21st issue edition (2026), arXiv:1401.0300v48. The registry describes the collection as edited by E. I. Khukhro and V. D. Mazurov. These bibliographic details are supplied attribution and have not been independently checked.
\item Comment of 2013 in the supplied problem statement: A. Lucchini, M. Morigi, P. Shumyatsky, \emph{Forum Math.} \textbf{24} (2012), 875--887. The comment attributes to this reference the existence of a two-generator subgroup with the same set of prime divisors as the ambient finite group. No article title or link is supplied.
\end{itemize}
\end{document}
